\chapter{Vaginal Bleeding}

\section*{Definition}
Vaginal bleeding is bleeding from the vagina that may be normal (menstruation) or abnormal. Abnormal vaginal bleeding includes bleeding between periods, postmenopausal bleeding, heavy menstrual bleeding, or bleeding after intercourse.

\section*{What Happens in Your Body}
Normal menstrual bleeding results from the shedding of the endometrial lining under the influence of declining progesterone. Abnormal bleeding can arise from hormonal imbalances (anovulation), structural abnormalities (fibroids, polyps), pregnancy complications, infection, or malignancy. Postmenopausal bleeding is caused by atrophic endometrium, hormone therapy, endometrial hyperplasia, or endometrial cancer.

\section*{Causes}
\begin{itemize}
\item Pregnancy-related: implantation bleeding, ectopic pregnancy, miscarriage, placental abruption
\item Hormonal: anovulation, perimenopause, thyroid disorders, hyperprolactinemia
\item Structural: uterine fibroids, endometrial polyps, cervical polyps, adenomyosis
\item Infectious: cervicitis, endometritis, pelvic inflammatory disease
\item Malignancy: cervical cancer, endometrial cancer, vaginal cancer
\item Medications: anticoagulants, tamoxifen, hormonal contraceptives
\item Trauma: sexual assault, foreign body, vaginal laceration
\item Postmenopausal bleeding: atrophic vaginitis, endometrial atrophy, endometrial hyperplasia or carcinoma
\end{itemize}

\section*{When to See a Doctor}
\begin{itemize}
\item Bleeding between menstrual periods
\item Heavy menstrual bleeding (soaking through pads/tampons hourly)
\item Bleeding after intercourse (postcoital bleeding)
\item Bleeding after menopause
\item Bleeding during pregnancy
\item Bleeding accompanied by pelvic pain or cramping
\item Clots larger than 1 inch in diameter
\item Signs of anemia: fatigue, lightheadedness, pale skin
\end{itemize}

\section*{Related Symptoms}
\begin{itemize}
\item Heavy menstrual bleeding (menorrhagia)
\item Irregular periods
\item Postmenopausal bleeding
\item Pelvic pain
\item Anemia
\item Uterine fibroids
\item Endometrial polyps
\item Cervical cancer
\end{itemize}

\section*{Self-Care / Home Management}
\begin{itemize}
\item Keep a menstrual diary to track bleeding patterns
\item Use pads or tampons and change frequently
\item Take over-the-counter pain relievers if needed
\item Maintain iron-rich diet if bleeding is heavy
\item Avoid aspirin (increases bleeding risk)
\item Seek medical attention for severe or persistent bleeding
\end{itemize}

\section*{How Your Doctor Will Evaluate You}
\begin{itemize}
\item Thorough history: menstrual pattern, pregnancy status, medications
\item Pregnancy test
\item Pelvic examination with speculum and bimanual exam
\item Complete blood count for anemia
\item Coagulation profile if bleeding disorder suspected
\item Pelvic ultrasound (transvaginal) for structural evaluation
\item Endometrial biopsy for postmenopausal bleeding or abnormal uterine bleeding
\item STI testing
\end{itemize}

\section*{Treatment Options}
\begin{itemize}
\item Treat underlying cause
\item Hormonal therapy: combined oral contraceptives, progestins, GnRH agonists
\item NSAIDs for heavy menstrual bleeding
\item Tranexamic acid to reduce bleeding
\item Surgical: hysteroscopy with polypectomy/myomectomy, endometrial ablation, hysterectomy
\item Dilation and curettage (D\&C) for acute heavy bleeding
\item Iron supplementation for anemia
\end{itemize}

\section*{References}
\begin{thebibliography}{99}
\bibitem{ref1} American College of Obstetricians and Gynecologists. "Diagnosis of abnormal uterine bleeding in reproductive-aged women." \textit{ACOG Practice Bulletin No. 128}, 2012;120(1):197--206.
\bibitem{ref2} Munro MG, Critchley HOD, Broder MS, et al. "FIGO classification system (PALM-COEIN) for causes of abnormal uterine bleeding in nongravid women." \textit{International Journal of Gynecology \& Obstetrics}, 2011;113(1):3--13.
\bibitem{ref3} Whitaker L, Critchley HOD. "Abnormal uterine bleeding." \textit{Best Practice \& Research Clinical Obstetrics \& Gynaecology}, 2016;34:47--57.
\bibitem{ref4} Goldstein SR, Lumsden MA. "Abnormal uterine bleeding in perimenopausal women." \textit{New England Journal of Medicine}, 2017;376(7):687--689.
\bibitem{ref5} Rahn DD, Ward RM, Sanses TVD, et al. "Vaginal bleeding in postmenopausal women: A systematic review." \textit{Obstetrics \& Gynecology}, 2018;131(4):e123--e133.
\bibitem{ref6} Sweet MG, Schmidt-Dalton TA, Weiss PM, et al. "Evaluation and management of abnormal uterine bleeding in premenopausal women." \textit{American Family Physician}, 2012;85(1):35--43.
\end{thebibliography}

\clearpage
